\documentclass{amsart}
% For dealing with references we use the comment environment
\usepackage{verbatim}
\newenvironment{reference}{\comment}{\endcomment}
%\newenvironment{reference}{}{}
\newenvironment{slogan}{\comment}{\endcomment}
\newenvironment{history}{\comment}{\endcomment}

% For commutative diagrams we use Xy-pic
\usepackage[all]{xy}

% We use 2cell for 2-commutative diagrams.
\xyoption{2cell}
\UseAllTwocells

% We use multicol for the list of chapters between chapters
\usepackage{multicol}

% This is generally recommended for better output
\usepackage{lmodern}
\usepackage[T1]{fontenc}

% For cross-file-references

% Package for hypertext links:
\usepackage{hyperref}

% For any local file, say "hello.tex" you want to link to please

% Theorem environments.
%
\theoremstyle{plain}
\newtheorem{theorem}[subsection]{Theorem}
\newtheorem{proposition}[subsection]{Proposition}
\newtheorem{lemma}[subsection]{Lemma}

\theoremstyle{definition}
\newtheorem{definition}[subsection]{Definition}
\newtheorem{example}[subsection]{Example}
\newtheorem{exercise}[subsection]{Exercise}
\newtheorem{situation}[subsection]{Situation}

\theoremstyle{remark}
\newtheorem{remark}[subsection]{Remark}
\newtheorem{remarks}[subsection]{Remarks}

\numberwithin{equation}{subsection}

% Macros
%
\def\lim{\mathop{\mathrm{lim}}\nolimits}
\def\colim{\mathop{\mathrm{colim}}\nolimits}
\def\Spec{\mathop{\mathrm{Spec}}}
\def\Hom{\mathop{\mathrm{Hom}}\nolimits}
\def\Ext{\mathop{\mathrm{Ext}}\nolimits}
\def\SheafHom{\mathop{\mathcal{H}\!\mathit{om}}\nolimits}
\def\SheafExt{\mathop{\mathcal{E}\!\mathit{xt}}\nolimits}
\def\Sch{\mathit{Sch}}
\def\Mor{\mathop{\mathrm{Mor}}\nolimits}
\def\Ob{\mathop{\mathrm{Ob}}\nolimits}
\def\Sh{\mathop{\mathit{Sh}}\nolimits}
\def\NL{\mathop{N\!L}\nolimits}
\def\CH{\mathop{\mathrm{CH}}\nolimits}
\def\proetale{{pro\text{-}\acute{e}tale}}
\def\etale{{\acute{e}tale}}
\def\QCoh{\mathit{QCoh}}
\def\Ker{\mathop{\mathrm{Ker}}}
\def\Im{\mathop{\mathrm{Im}}}
\def\Coker{\mathop{\mathrm{Coker}}}
\def\Coim{\mathop{\mathrm{Coim}}}

% Boxtimes
%
\DeclareMathSymbol{\boxtimes}{\mathbin}{AMSa}{"02}

%
% Macros for moduli stacks/spaces
%
\def\QCohstack{\mathcal{QC}\!\mathit{oh}}
\def\Cohstack{\mathcal{C}\!\mathit{oh}}
\def\Spacesstack{\mathcal{S}\!\mathit{paces}}
\def\Quotfunctor{\mathrm{Quot}}
\def\Hilbfunctor{\mathrm{Hilb}}
\def\Curvesstack{\mathcal{C}\!\mathit{urves}}
\def\Polarizedstack{\mathcal{P}\!\mathit{olarized}}
\def\Complexesstack{\mathcal{C}\!\mathit{omplexes}}
% \Pic is the operator that assigns to X its picard group, usage \Pic(X)
% \Picardstack_{X/B} denotes the Picard stack of X over B
% \Picardfunctor_{X/B} denotes the Picard functor of X over B
\def\Pic{\mathop{\mathrm{Pic}}\nolimits}
\def\Picardstack{\mathcal{P}\!\mathit{ic}}
\def\Picardfunctor{\mathrm{Pic}}
\def\Deformationcategory{\mathcal{D}\!\mathit{ef}}

\setlength{\emergencystretch}{3em}
\begin{document}
\title[Stacks Project, Tag 0BJ6]{729 --- Generic smoothness over unramified DVR extensions with separable residue fields is resolved by Neron blowups}
\maketitle
\noindent Copyright (C) 2005--2025 Johan de Jong. Stacks Project authors. Source: \href{https://stacks.math.columbia.edu/tag/0BJ6}{Tag 0BJ6}.

\noindent Editorial difficulty note (not author text). Difficulty H3: The proof assigns a torsion-length defect to a differential presentation and proves strict decrease after each Neron blowup by tracking free lattices and divisibility of minimal generators. This terminating desingularization argument is research-level arithmetic geometry..

\noindent Editorial provenance note (not author text). History: excerpt from revision \texttt{a0444\allowbreak 6e57e\allowbreak c1fbc\allowbreak 252a8\allowbreak 71afc\allowbreak ec775\allowbreak 2fb28\allowbreak 07b14}, smoothing.tex, lines 1040--1176. Reference presentation was adapted; original-author context and core support blocks were retained or added. The mathematical wording is unchanged. Licensed under GNU FDL 1.2 or later; no invariant sections or cover texts. See ../licenses/COPYING. Context blocks reproduce author definitions and supporting results; they are not additional counted problems.

\begin{situation}
\label{situation-neron}
Here $R \subset \Lambda$ is an extension of discrete valuation rings
with ramification index $1$ (More on Algebra, Definition
\ref{more-algebra-definition-extension-discrete-valuation-rings}).
We assume given a factorization
$$
R \to A \xrightarrow{\varphi} \Lambda
$$
with $R \to A$ flat and of finite type. Let $\mathfrak q = \Ker(\varphi)$
and $\mathfrak p = \varphi^{-1}(\mathfrak m_\Lambda)$.
\end{situation}

\begin{definition}
\label{algebra-definition-relative-dimension}
Suppose that $R \to S$ is of finite type, and let
$\mathfrak q \subset S$ be a prime lying over a prime
$\mathfrak p$ of $R$.
We define the {\it relative dimension
of $S/R$ at $\mathfrak q$}, denoted
$\dim_{\mathfrak q}(S/R)$, to be the dimension
of $\Spec(S \otimes_R \kappa(\mathfrak p))$
at the point corresponding to $\mathfrak q$. We let
$\dim(S/R)$ be the supremum of $\dim_{\mathfrak q}(S/R)$
over all $\mathfrak q$. This is called the
{\it relative dimension of} $S/R$.
\end{definition}

\begin{definition}
\label{algebra-definition-smooth}
A ring map $R \to S$ is {\it smooth} if it is of finite presentation
and the naive cotangent complex $\NL_{S/R}$ is quasi-isomorphic to a
finite projective $S$-module placed in degree $0$: this means
that $H_1(\NL_{S/R}) = 0$ and that $\Omega_{S/R}$ is a finite projective
$S$-module.
\end{definition}

\begin{definition}
\label{algebra-definition-separable-field-extension}
Let $K/k$ be a field extension.
\begin{enumerate}
\item We say $K$ is {\it separably generated over $k$} if there exists
a transcendence basis $\{x_i; i \in I\}$ of $K/k$ such that the extension
$K/k(x_i; i \in I)$ is a separable algebraic extension.
\item We say $K$ is {\it separable over $k$} if for every subextension
$k \subset K' \subset K$ with $K'$ finitely generated
over $k$, the extension $K'/k$ is separably generated.
\end{enumerate}
\end{definition}

\begin{definition}
\label{more-algebra-definition-extension-discrete-valuation-rings}
We say that $A \to B$ or $A \subset B$ is an
{\it extension of discrete valuation rings} if $A$ and $B$ are
discrete valuation rings and $A \to B$ is injective and local.
In particular, if $\pi_A$ and $\pi_B$ are uniformizers of
$A$ and $B$, then $\pi_A = u \pi_B^e$ for some $e \geq 1$ and unit
$u$ of $B$. The integer $e$ does not depend on the choice of
the uniformizers as it is also the unique integer $\geq 1$ such that
$$
\mathfrak m_A B = \mathfrak m_B^e
$$
The integer $e$ is called the {\it ramification index} of $B$ over $A$.
We say that $B$ is {\it weakly unramified} over $A$ if $e = 1$.
If the extension of residue fields
$\kappa_A = A/\mathfrak m_A \subset \kappa_B = B/\mathfrak m_B$
is finite, then we set $f = [\kappa_B : \kappa_A]$ and we
call it the {\it residual degree} or {\it residue degree}
of the extension $A \subset B$.
\end{definition}

\begin{lemma}
\label{lemma-neron-blowup-smooth}
In Situation \ref{situation-neron} assume that $R \to A$ is smooth
at $\mathfrak p$ and that $R/\pi R \subset \Lambda/\pi \Lambda$
is a separable field extension. Then $R \to A'$ is smooth at
$\mathfrak p'$ and there is a short exact sequence
$$
0 \to
\Omega_{A/R} \otimes_A A'_{\mathfrak p'} \to
\Omega_{A'/R, \mathfrak p'} \to
(A'/\pi A')_{\mathfrak p'}^{\oplus c} \to 0
$$
where $c = \dim((A/\pi A)_\mathfrak p)$.
\end{lemma}

\begin{proof}
By Lemma \href{https://stacks.math.columbia.edu/tag/0BJ3}{lemma neron functorial (Tag 0BJ3)} we may replace $A$ by a localization
at an element not in $\mathfrak p$; we will use this without further mention.
Write $\kappa = R/\pi R$. Since smoothness is stable under base change
(Algebra, Lemma \href{https://stacks.math.columbia.edu/tag/00T4}{lemma base change smooth (Tag 00T4)})
we see that $A/\pi A$ is smooth over $\kappa$ at $\mathfrak p$.
Hence $(A/\pi A)_\mathfrak p$ is a regular local ring
(Algebra, Lemma \href{https://stacks.math.columbia.edu/tag/00TT}{lemma characterize smooth over field (Tag 00TT)}).
Choose $g_1, \ldots, g_c \in \mathfrak p$ which map to
a regular system of parameters in $(A/\pi A)_\mathfrak p$.
Then we see that $\mathfrak p = (\pi, g_1, \ldots, g_c)$
after possibly replacing $A$ by a localization.
Note that $\pi, g_1, \ldots, g_c$ is a regular sequence
in $A_\mathfrak p$ (first $\pi$ is a nonzerodivisor and
then Algebra, Lemma \href{https://stacks.math.columbia.edu/tag/00NQ}{lemma regular ring CM (Tag 00NQ)}
for the rest of the sequence).
After replacing $A$ by a localization we may assume that
$\pi, g_1, \ldots, g_c$ is a regular sequence in $A$
(Algebra, Lemma \href{https://stacks.math.columbia.edu/tag/061L}{lemma regular sequence in neighbourhood (Tag 061L)}).
It follows that
$$
A' = A[y_1, \ldots, y_c]/(\pi y_1 - g_1, \ldots, \pi y_c - g_c) =
A[y_1, \ldots, y_c]/I
$$
by More on Algebra, Lemma \href{https://stacks.math.columbia.edu/tag/0BIQ}{lemma blowup regular sequence (Tag 0BIQ)}.
In the following we will use the definition of smoothness using the
naive cotangent complex (Algebra, Definition \ref{algebra-definition-smooth})
and the criterion of Algebra, Lemma \href{https://stacks.math.columbia.edu/tag/07BU}{lemma smooth at point (Tag 07BU)}
without further mention.
The exact sequence of Algebra, Lemma \href{https://stacks.math.columbia.edu/tag/00S2}{lemma exact sequence NL (Tag 00S2)}
for $R \to A[y_1, \ldots, y_c] \to A'$ looks like this
$$
0 \to H_1(\NL_{A'/R}) \to I/I^2 \to
\Omega_{A/R} \otimes_A A' \oplus
\bigoplus\nolimits_{i = 1, \ldots, c} A' \text{d}y_i \to
\Omega_{A'/R} \to 0
$$
where the class of $\pi y_i - g_i$ in $I/I^2$ is mapped
to $- \text{d}g_i + \pi \text{d}y_i$ in the next term.
Here we have used Algebra, Lemma \href{https://stacks.math.columbia.edu/tag/07BP}{lemma NL surjection (Tag 07BP)}
to compute $\NL_{A'/A[y_1, \ldots, y_c]}$ and we have used that
$R \to A[y_1, \ldots, y_c]$ is smooth, so
$H_1(\NL_{A[y_1, \ldots, y_c]/R}) = 0$ and
$\Omega_{A[y_1, \ldots, y_c]/R}$ is a finite projective (a fortiori flat)
$A[y_1, \ldots, y_c]$-module which is in fact the direct sum
of $\Omega_{A/R} \otimes_A A[y_1, \ldots, y_c]$ and a free
module with basis $\text{d}y_i$. To finish the proof it
suffices to show that $\text{d}g_1, \ldots, \text{d}g_c$
forms part of a basis for the finite free module $\Omega_{A/R, \mathfrak p}$.
Namely, this will show $(I/I^2)_\mathfrak p$ is free on $\pi y_i - g_i$,
the localization at $\mathfrak p$ of the middle map in the sequence is
injective, so $H_1(\NL_{A'/R})_\mathfrak p = 0$, and
that the cokernel $\Omega_{A'/R, \mathfrak p}$ is finite free.
To do this it suffices to show
that the images of $\text{d}g_i$ are $\kappa(\mathfrak p)$-linearly
independent in
$\Omega_{A/R, \mathfrak p}/\pi = \Omega_{(A/\pi A)/\kappa, \mathfrak p}$
(equality by Algebra, Lemma \href{https://stacks.math.columbia.edu/tag/00RV}{lemma differentials base change (Tag 00RV)}).
Since $\kappa \subset \kappa(\mathfrak p) \subset \Lambda/\pi \Lambda$
we see that $\kappa(\mathfrak p)$ is separable over $\kappa$
(Algebra, Definition \ref{algebra-definition-separable-field-extension}).
The desired linear independence now follows from
Algebra, Lemma \href{https://stacks.math.columbia.edu/tag/00TU}{lemma computation differential (Tag 00TU)}.
\end{proof}


\begin{lemma}
\label{lemma-neron-when-smooth}
In Situation \ref{situation-neron} assume that $R \to A$ is smooth
at $\mathfrak q$ and that we have a surjection of $R$-algebras
$B \to A$ with kernel $I$. Assume $R \to B$ smooth at
$\mathfrak p_B = (B \to A)^{-1}\mathfrak p$. If the cokernel of
$$
I/I^2 \otimes_A \Lambda \to \Omega_{B/R} \otimes_B \Lambda
$$
is a free $\Lambda$-module, then $R \to A$ is smooth at $\mathfrak p$.
\end{lemma}

\begin{proof}
The cokernel of the map $I/I^2 \to \Omega_{B/R} \otimes_B A$
is $\Omega_{A/R}$, see Algebra, Lemma \href{https://stacks.math.columbia.edu/tag/00RU}{lemma differential seq (Tag 00RU)}.
Let $d = \dim_\mathfrak q(A/R)$ be the relative dimension of $R \to A$
at $\mathfrak q$, i.e., the dimension of $\Spec(A[1/\pi])$ at $\mathfrak q$.
See Algebra, Definition \ref{algebra-definition-relative-dimension}.
Then $\Omega_{A/R, \mathfrak q}$ is free over $A_\mathfrak q$ of rank $d$
(Algebra, Lemma \href{https://stacks.math.columbia.edu/tag/00TT}{lemma characterize smooth over field (Tag 00TT)}).
Thus if the hypothesis of the lemma holds,
then $\Omega_{A/R} \otimes_A \Lambda$ is free of rank $d$.
It follows that $\Omega_{A/R} \otimes_A \kappa(\mathfrak p)$
has dimension $d$ (as it is true upon tensoring with $\Lambda/\pi \Lambda$).
Since $R \to A$ is flat and since $\mathfrak p$
is a specialization of $\mathfrak q$, we see that
$\dim_\mathfrak p(A/R) \geq d$ by Algebra, Lemma
\href{https://stacks.math.columbia.edu/tag/00QH}{lemma dimension fibres bounded open upstairs (Tag 00QH)}.
Then it follows that $R \to A$ is smooth at $\mathfrak p$ by
Algebra, Lemmas \href{https://stacks.math.columbia.edu/tag/00TF}{lemma flat fibre smooth (Tag 00TF)} and
\href{https://stacks.math.columbia.edu/tag/00TT}{lemma characterize smooth over field (Tag 00TT)}.
\end{proof}


\noindent
In Situation \ref{situation-neron} let $\pi \in R$ be a uniformizer.
Recall that flatness of $A$ over $R$ signifies that $\pi$
is a nonzerodivisor on $A$
(More on Algebra, Lemma
\href{https://stacks.math.columbia.edu/tag/0539}{lemma valuation ring torsion free flat (Tag 0539)}).
By our assumption on $R \subset \Lambda$ we see that $\pi$ maps
to a uniformizer of $\Lambda$. Since $\pi \in \mathfrak p$ we can consider
N\'eron's affine blowup algebra (see
Algebra, Section \href{https://stacks.math.columbia.edu/tag/052P}{section blow up (Tag 052P)})
$$
\varphi' :
A' = A[\textstyle{\frac{\mathfrak p}{\pi}}]
\longrightarrow
\Lambda
$$
which comes endowed with an induced map to $\Lambda$ sending
$a/\pi^n$, $a \in \mathfrak p^n$ to $\pi^{-n}\varphi(a)$ in $\Lambda$.
We will denote $\mathfrak q' \subset \mathfrak p' \subset A'$
the corresponding prime ideals of $A'$. Observe that the isomorphism
class of $A'$ does not depend on our choice of uniformizer.
Repeating the construction we obtain a sequence
$$
A \to A' \to A'' \to \ldots \to \Lambda
$$


\begin{lemma}
\label{lemma-neron-desingularization}
In Situation \ref{situation-neron}
assume that $R \to A$ is smooth at $\mathfrak q$
and that $R/\pi R \subset \Lambda/\pi \Lambda$ is a separable
extension of fields. Then after a finite number of affine N\'eron
blowups the algebra $A$ becomes smooth over $R$ at $\mathfrak p$.
\end{lemma}

\begin{proof}
We choose an $R$-algebra $B$ and a surjection $B \to A$. Set
$\mathfrak p_B = (B \to A)^{-1}(\mathfrak p)$ and denote $r$
the relative dimension of $R \to B$ at $\mathfrak p_B$. We choose $B$
such that $R \to B$ is smooth at $\mathfrak p_B$.
For example we can take $B$ to be a polynomial algebra in $r$ variables
over $R$. Consider the complex
$$
I/I^2 \otimes_A \Lambda \longrightarrow \Omega_{B/R} \otimes_B \Lambda
$$
of Lemma \ref{lemma-neron-when-smooth}. By the structure of finite modules
over $\Lambda$ (More on Algebra, Lemma \href{https://stacks.math.columbia.edu/tag/0ASV}{lemma modules PID (Tag 0ASV)})
we see that the cokernel looks like
$$
\Lambda^{\oplus d} \oplus
\bigoplus\nolimits_{i = 1, \ldots, n} \Lambda/\pi^{e_i} \Lambda
$$
for some $d \geq 0$, $n \geq 0$, and $e_i \geq 1$. Observe that $d$
is the relative dimension of $A/R$ at $\mathfrak q$
(Algebra, Lemma \href{https://stacks.math.columbia.edu/tag/00TT}{lemma characterize smooth over field (Tag 00TT)}).
If the defect $e = \sum_{i = 1, \ldots, n} e_i$ is zero, then we are done by
Lemma \ref{lemma-neron-when-smooth}.

\medskip\noindent
Next, we consider what happens when we perform the N\'eron blowup.
Recall that $A'$ is the quotient of $B'/IB'$ by its $\pi$-power
torsion (Lemma \href{https://stacks.math.columbia.edu/tag/0BJ3}{lemma neron functorial (Tag 0BJ3)}) and that $R \to B'$ is smooth at
$\mathfrak p_{B'}$ (Lemma \ref{lemma-neron-blowup-smooth}).
Thus after blowup we have exactly the same setup. Picture
$$
\xymatrix{
0 \ar[r] & I' \ar[r] & B' \ar[r] & A' \ar[r] & 0 \\
0 \ar[r] & I \ar[u] \ar[r] & B \ar[u] \ar[r] & A \ar[r] \ar[u] & 0
}
$$
Since $I \subset \mathfrak p_B$, we see that $I \to I'$
factors through $\pi I'$. Looking at the induced map of
complexes we get
$$
\xymatrix{
I'/(I')^2 \otimes_{A'} \Lambda \ar[r] &
\Omega_{B'/R} \otimes_{B'} \Lambda \ar@{=}[r] & M' \\
I/I^2 \otimes_A \Lambda \ar[r] \ar[u] &
\Omega_{B/R} \otimes_B \Lambda \ar[u] \ar@{=}[r] & M
}
$$
Then $M \subset M'$ are finite free $\Lambda$-modules with quotient
$M'/M$ annihilated by $\pi$, see Lemma \ref{lemma-neron-blowup-smooth}.
Let $N \subset M$ and $N' \subset M'$ be the images of the horizontal
maps and denote $Q = M/N$ and $Q' = M'/N'$.
We obtain a commutative diagram
$$
\xymatrix{
0 \ar[r] &
N' \ar[r] &
M' \ar[r] &
Q' \ar[r] &
0 \\
0 \ar[r] &
N \ar[r] \ar[u] &
M \ar[r] \ar[u] &
Q \ar[r] \ar[u] &
0
}
$$
Then $N \subset N'$ are free $\Lambda$-modules of rank $r - d$.
Since $I$ maps into $\pi I'$ we see that $N \subset \pi N'$.

\medskip\noindent
Let $K = \Lambda_\pi$ be the fraction field of $\Lambda$.
We have a commutative diagram
$$
\xymatrix{
0 \ar[r] &
N' \ar[r] &
N'_K \cap M' \ar[r] &
Q'_{tor} \ar[r] &
0 \\
0 \ar[r] &
N \ar[r] \ar[u] &
N_K \cap M \ar[r] \ar[u] &
Q_{tor} \ar[r] \ar[u] &
0
}
$$
whose rows are short exact sequences. This shows that the change in defect
is given by
$$
e - e' =
\text{length}(Q_{tor}) - \text{length}(Q'_{tor})
=
\text{length}(N'/N) - \text{length}(N'_K \cap M' / N_K \cap M)
$$
Since $M'/M$ is annihilated by $\pi$, so is $N'_K \cap M' / N_K \cap M$,
and its length is at most $\dim_K(N_K)$.
Since $N \subset \pi N'$ we get $\text{length}(N'/N) \ge \dim_K(N_K)$,
with equality if and only if $N = \pi N'$.

\medskip\noindent
To finish the proof we have to show that $N$ is strictly smaller
than $\pi N'$ when $A$ is not smooth at $\mathfrak p$;
this is the key computation one has to do in N\'eron's argument.
To do this, we consider the exact sequence
$$
I/I^2 \otimes_B \kappa(\mathfrak p_B)
\to \Omega_{B/R} \otimes_B \kappa(\mathfrak p_B)
\to \Omega_{A/R} \otimes_A \kappa(\mathfrak p) \to 0
$$
(follows from Algebra, Lemma \href{https://stacks.math.columbia.edu/tag/00RU}{lemma differential seq (Tag 00RU)}).
Since $R \to A$ is not smooth at $\mathfrak p$ we see that the dimension $s$ of
$\Omega_{A/R} \otimes_A \kappa(\mathfrak p)$
is bigger than $d$. On the other hand
the first arrow factors through the injective map
$$
\mathfrak p B_\mathfrak p/\mathfrak p^2 B_\mathfrak p
\to \Omega_{B/R} \otimes_B \kappa(\mathfrak p_B)
$$
of Algebra, Lemma \href{https://stacks.math.columbia.edu/tag/00TU}{lemma computation differential (Tag 00TU)};
note that $\kappa(\mathfrak p)$ is separable over $k$
by our assumption on $R/\pi R \subset \Lambda/\pi \Lambda$.
Hence we conclude that we can find generators
$g_1, \ldots, g_t \in I$ such that $g_j \in \mathfrak p^2$
for $j > r - s$. Then the images of $g_j$ in $A'$ are in $\pi^2 I'$
for $j > r - s$. Since $r - s < r - d$
we find that at least one of the minimal generators
of $N$ becomes divisible by $\pi^2$ in $N'$.
Thus we see that $e$ decreases by at least $1$ and we win.
\end{proof}
\end{document}
